\subsection{Degree reduction for the UGC independent-set source}
\label{app:bounded-degree-ugc-source}

We derive Theorem~\ref{thm:bounded-degree-kr-source} from the
weighted density theorem of Austrin, Khot, and Safra. Their later
bounded-degree reduction uses random sampling
\cite[Section~4, Steps~2--3]{AKS2011}. We instead obtain bounded
degree deterministically by spectral sparsification. We state the
cited density theorem first, then derive the loop removal, vertex
unweighting, and degree bounds needed here.

\subsubsection{The AKS density theorem}

A bipartite \emph{Unique Games instance}
$\Lambda=(U,V,E_\Lambda,\Pi,[\ell])$ consists of a bipartite
multigraph with left vertices $U$, right vertices $V$, and a label
set $[\ell]$. Each edge $e=uv$ has a permutation $\pi_e\in\Pi$
and requires the label at $v$ to be $\pi_e$ applied to the label at
$u$. Write $\OPT(\Lambda)$ for the maximum fraction
of satisfied constraints, counting parallel constraints separately.
Write $\operatorname{VOPT}(\Lambda)$ for the largest fraction
$|V_0|/|V|$ for which a labeling satisfies every constraint incident
to $V_0\subseteq V$. The instance is \emph{regular} if all left
vertices have the same degree and all right vertices have the same
degree. The two degrees may differ.

\begin{theorem}[Regular Unique Games source
{\cite[Conjecture~2.4]{AKS2011} and \cite{KhotRegev}}]
\label{thm:regular-ug-source}
For every fixed $\gamma>0$, there is a fixed alphabet size $\ell$
such that it is UGC-hard to distinguish regular bipartite Unique
Games instances $\Lambda$ with alphabet $[\ell]$ satisfying the
following cases.
\begin{description}
\item[YES:] $\operatorname{VOPT}(\Lambda)\geq1-\gamma$.
 There is a set $V_0\subseteq V$ of at least $(1-\gamma)|V|$ right
 vertices and one labeling satisfying every constraint incident
 to $V_0$.
\item[NO:] $\OPT(\Lambda)\leq\gamma$.
 Every labeling satisfies at most a $\gamma$ fraction of the
 constraints.
\end{description}
\end{theorem}

For a vertex- and edge-weighted graph $G$, let $a_v$ denote a
vertex weight and $w_e$ an edge weight, normalized so that each
of the two total weights is one. Put
$a(S)=\sum_{v\in S}a_v$ and let $e_G(S)$ be the total weight of
edges with both endpoints in $S$. We allow loops, counted once in
$e_G(S)$ and twice in the weighted degree $\deg_G(v)$.
For $\beta\in[0,1]$ and a real $\delta$, the graph is
\emph{$(\beta,\delta)$-dense} if $a(S)\geq\beta$ implies
$e_G(S)\geq\delta$.

Let $\Psi$ be the standard normal cumulative distribution function.
For $-1<\rho<1$, let $X,Y$ be jointly normal with mean zero,
variance one, and covariance $\rho$, and define
\[
 \Gamma_\rho(s)=
 \Pr[X\leq\Psi^{-1}(s),\ Y\leq\Psi^{-1}(s)]
 \qquad(0<s<1).
\]
Set $\Gamma_\rho(0)=0$ and $\Gamma_\rho(1)=1$.
Fix a bias $q\in(0,1/2)$ and an analytic error
$\varepsilon_{\mathrm{dens}}>0$, and put $\rho=-q/(1-q)$.

\Needspace{11\baselineskip}
\begin{theorem}[AKS density theorem {\cite[Theorem~3.1]{AKS2011}}]
\label{thm:aks-density-source}
For all sufficiently small $\gamma>0$, an algorithm maps a Unique
Games instance $\Lambda$ to a vertex- and edge-weighted graph $G$
with the following properties.
\begin{description}
\item[YES:] If $\operatorname{VOPT}(\Lambda)\geq1-\gamma$, then
 $G$ has an independent set of vertex weight at least $q-\gamma$.
\item[NO:] If $\OPT(\Lambda)\leq\gamma$ and
 $\Lambda$ is regular, then $G$ is
 $(\beta,\Gamma_\rho(\beta)-\varepsilon_{\mathrm{dens}})$-dense
 for every $\beta\in[0,1]$.
\end{description}
Vertex weights are proportional to weighted degrees.
The running time is polynomial in $|U|,|V|,|E_\Lambda|$ and
exponential in $\ell$.
\end{theorem}

For a regular input, the proportionality gives
$\deg_G(v)=2a_v$: the weighted degrees sum to twice the total edge
weight, whereas the vertex weights sum to one. Equivalently, each
endpoint marginal of the AKS edge distribution is $a_v$, with a
loop contributing twice. Thus the NO conclusion, applied with
$\beta=a(S)$, is
\[
 e_G(S)\geq\Gamma_\rho(a(S))-\varepsilon_{\mathrm{dens}}
 \qquad(S\subseteq V(G)).
\]

We will use positivity, monotonicity, and convexity of $\Gamma_\rho$.
For completeness, these follow directly from its definition. The
joint normal density is positive, so $\Gamma_\rho(s)>0$ for
$0<s<1$. Differentiating the bivariate distribution function along
its diagonal, using symmetry of $X,Y$ and the conditional normal
law, gives
\[
 \Gamma_\rho'(s)=
 2\Psi\!\left(\sqrt{\frac{1-\rho}{1+\rho}}\,\Psi^{-1}(s)\right).
\]
The right-hand side is nonnegative and increasing in $s$.
Consequently $\Gamma_\rho$ is nondecreasing and convex, including
at the endpoints by continuity.

\subsubsection{Removing loops and vertex weights}

Fix a rational bias $q$ from now on. The AKS sampling construction
can produce loops. We remove them before replacing vertices by
independent copies, so that the density argument can use
independent choices of distinct original vertices. Let
$\varepsilon_{\mathrm{loop}}>0$ be a fixed loop-loss budget.
The following alphabet padding makes the total loop weight smaller
than this budget.

\begin{lemma}[Removing loops by alphabet padding]\label{lem:aks-loop-removal}
For every fixed $\varepsilon_{\mathrm{loop}}>0$, a fixed alphabet padding followed by the
AKS construction and deletion of all loops gives a loopless graph $G^-$
with vertex weights $a_v$ of total one, total edge weight at most
one, and weighted degrees at most $2a_v$.
Its YES independent-set weight is at least $q-\gamma$, and in the
NO case
\begin{equation}\label{eq:aks-loop-free-density}
 e_{G^-}(S)\ge\Gamma_\rho(a(S))-\varepsilon_{\mathrm{dens}}-\varepsilon_{\mathrm{loop}}
 \qquad(S\subseteq V(G^-)).
\end{equation}
The reduction is deterministic and polynomial for fixed parameters.
\end{lemma}

\begin{proof}
To control possible loops even when the Unique Games instance has
parallel constraints with different permutations, first pad its
alphabet from $[\ell]$ to $[\ell]\times[t]$, replacing every
constraint $\pi$ by $(i,j)\mapsto(\pi(i),j)$.
Projecting any satisfying padded label pair to its first coordinate
shows that padding cannot increase the NO optimum. A YES labeling
extends by giving every labeled vertex the same second coordinate.
Regularity and the YES fraction are unchanged.

The AKS edge distribution~\cite[proof of Theorem~3.1]{AKS2011}
first chooses two constraints incident with the same left vertex,
then samples independent coordinate pairs $(x_j,y_j)$ with
probabilities $\Pr[x_j=0,y_j=0]=1-2q$,
$\Pr[x_j=0,y_j=1]=\Pr[x_j=1,y_j=0]=q$, and
$\Pr[x_j=1,y_j=1]=0$. It then permutes the endpoint strings by
the two constraints.

The coordinates $[\ell]\times\{j\}$, for a fixed $j\in[t]$,
form one coordinate block. Put $c=q^2+(1-q)^2<1$. Conditional on any choice of the two
constraints, a loop requires equality of the two permuted strings
in every one of the $t$ coordinate blocks. In each block, full
string equality implies equality at a fixed output coordinate. If
the two permutations select the same input coordinate there, equality
has probability $1-2q$. If they select different input coordinates,
the bits are independent with marginal probability $q$, so equality
has probability $c=q^2+(1-q)^2$. Both bounds are at most $c$.
Indeed, $c-(1-2q)=2q^2>0$.
The coordinate blocks are independent, so the probability of full
string equality, and hence the total loop weight, is at most $c^t$.
Choose a fixed $t$ with $c^t<\varepsilon_{\mathrm{loop}}$, and delete all loops.
This loses at most $\varepsilon_{\mathrm{loop}}$ internal edge weight in any set,
preserves the YES independent set, and does not increase weighted
degrees. Alphabet padding preserves the source promise and
regularity, so the construction above has the stated bounds.
All operations are deterministic, and $t$ is fixed.
\end{proof}

\subsubsection{Vertex unweighting and uniform degree bounds}

Write the padded alphabet size as $R$ and the fixed bias as
$q=a_0/a_1$, with positive integers $a_0<a_1$. For a binary
string $x\in\{0,1\}^R$, let $|x|$ be its
number of coordinates equal to one. The vertex weights in the AKS
construction, indexed by right vertices $v\in V$ and strings $x$,
have the form
\[
 a_{(v,x)}=
 \frac{a_0^{|x|}(a_1-a_0)^{R-|x|}}{|V|a_1^R},
 \qquad v\in V,\quad x\in\{0,1\}^R.
\]
Thus $N_{\mathrm{cl}}=|V|a_1^R$ is polynomial in the Unique Games
instance, since $a_1$ and $R$ are fixed. Replace each original vertex
$u$ by $m_u=N_{\mathrm{cl}}a_u$ independent copies, called
\emph{clones}, labeled
$(u,1),\ldots,(u,m_u)$. These copies form the \emph{clone block} of $u$.
If several source edges join the same pair of vertices, combine their
weights. For every resulting edge $uv$ of weight $w_{uv}$ in $G^-$,
join each clone of $u$ to each clone of $v$ by an edge of weight
$w_{uv}/(m_um_v)$. There are no other edges.
Denote this edge-weighted, vertex-unweighted
graph by $G_1$. It has $N_{\mathrm{cl}}$ vertices, total edge weight
at most one, and weighted degree at most $2/N_{\mathrm{cl}}$ at each vertex.

A clone of $u$ has weighted degree
$\deg_{G^-}(u)/m_u\le2a_u/(N_{\mathrm{cl}}a_u)=2/N_{\mathrm{cl}}$.
The YES independent set becomes one of the same normalized size by
selecting all its clone blocks.

\begin{lemma}[Density after vertex unweighting
{\cite[Section~4, Step~1]{AKS2011}}]
\label{lem:aks-clone-density}
If $G^-$ satisfies~\eqref{eq:aks-loop-free-density}, then its
clone graph $G_1$ satisfies, for every $S'\subseteq V(G_1)$,
\begin{equation}\label{eq:aks-clone-density}
 e_{G_1}(S')\ge\Gamma_\rho(|S'|/N_{\mathrm{cl}})
 -\varepsilon_{\mathrm{dens}}-\varepsilon_{\mathrm{loop}}.
\end{equation}
\end{lemma}

\begin{proof}
For a chosen set $S'\subseteq V(G_1)$, let $f_u$ be the fraction of
clones of $u$ in $S'$. The cited argument forms a random set
$S\subseteq V(G^-)$ by selecting each source vertex $u$ independently
with probability $f_u$. Since $a_u=m_u/N_{\mathrm{cl}}$,
each clone represents source weight $1/N_{\mathrm{cl}}$. The loop
removal in Lemma~\ref{lem:aks-loop-removal} makes the two endpoints
of each edge distinct, so their selection events are independent.
For an edge $uv$ of weight $w_{uv}$, each of the $m_um_v$ clone pairs
has weight $w_{uv}/(m_um_v)$. The set $S'$ contains $m_uf_u$ clones
of $u$ and $m_vf_v$ of $v$, so these clone pairs contribute
$(m_uf_u)(m_vf_v)w_{uv}/(m_um_v)=w_{uv}f_uf_v$. Thus the expected
source weight and expected internal edge weight are
\[
 \mathbb E[a(S)]=\sum_u a_uf_u=\frac{|S'|}{N_{\mathrm{cl}}},
 \qquad
 \mathbb E[e_{G^-}(S)]=\sum_{uv}w_{uv}f_uf_v=e_{G_1}(S').
\]
Applying the source density bound and then convexity of $\Gamma_\rho$
gives
\[
 \begin{aligned}
 e_{G_1}(S')
 &=\mathbb E[e_{G^-}(S)]\\
 &\geq\mathbb E[\Gamma_\rho(a(S))]
       -\varepsilon_{\mathrm{dens}}-\varepsilon_{\mathrm{loop}}\\
 &\geq\Gamma_\rho\!\left(\mathbb E[a(S)]\right)
       -\varepsilon_{\mathrm{dens}}-\varepsilon_{\mathrm{loop}}\\
 &=\Gamma_\rho(|S'|/N_{\mathrm{cl}})
       -\varepsilon_{\mathrm{dens}}-\varepsilon_{\mathrm{loop}}.
 \end{aligned}
\]
This is the claimed bound.
\end{proof}

For the normalized promise, fix a rational $0<\eta<1/10$.
Let $N$ denote the number of output vertices. The positive rational
constant $\zeta_\eta$ below may depend on $\eta$.

\begin{lemma}[Normalized density promise]
\label{lem:normalized-density-source}
There is a choice of $\zeta_\eta$ for which it is UGC-hard to
distinguish simple edge-weighted graphs on $N$
vertices, with positive rational edge weights, total edge weight
at most one, and weighted degree at most $2/N$, satisfying
\begin{description}
\item[YES:] There is an independent set of size at least $(1/2-\eta)N$.
\item[NO:] Every set of at least $\eta N$ vertices has internal edge
 weight at least $\zeta_\eta$.
\end{description}
The reduction is deterministic and polynomial time.
\end{lemma}

\begin{proof}
Set $q=1/2-\eta/4$ and $\rho=-q/(1-q)$. Choose a positive rational
$g$ with $4g<\Gamma_\rho(\eta)$, take
$\varepsilon_{\mathrm{dens}}\leq g$ and
$\varepsilon_{\mathrm{loop}}=g$, and choose $\gamma\leq\eta/4$
small enough for Theorem~\ref{thm:aks-density-source}.
Apply the regular source of Theorem~\ref{thm:regular-ug-source},
then the loop removal and clone construction above.
The YES independent-set density is at least
$q-\gamma\geq1/2-\eta/2\geq1/2-\eta$. For $|S|\geq\eta N$,
equation~\eqref{eq:aks-clone-density} and monotonicity give
\[
 e_{G_1}(S)\geq\Gamma_\rho(\eta)-2g>2g>g.
\]
Thus we may take $\zeta_\eta=g$.

The loop removal makes the clone graph simple after parallel edge
weights are combined. Its total edge weight and degree bounds were
proved above. Rational bias gives rational probabilities in the
finite AKS edge distribution. Enumerating the choices computes the
edge weights exactly. The clone weights are then rational as well.
For fixed $\eta$, all alphabet sizes, denominators of the bias,
and padding factors are constants. The displayed clone count is
therefore polynomial, and all operations are deterministic.
\end{proof}

\subsubsection{Deterministic sparsification of a dense-set promise}

To obtain bounded ordinary degree, we sparsify the weighted graph
and discard edges of very small weight. There will be few enough
edges that discarding them loses little total weight. Every remaining
edge then has a lower weight bound, so a bound on weighted degree
limits the number of incident edges. We use spectral sparsification
to preserve the internal weight of large vertex sets throughout this
process.

For a loopless weighted graph $G$,
let $w_{uv}$ be the total weight of edges between distinct vertices
$u,v$. Its \emph{Laplacian} is the symmetric matrix with
$L_G(u,u)=\deg_G(u)$ and $L_G(u,v)=-w_{uv}$ for $u\ne v$.
Expanding the quadratic form gives
\[
 x^{\mathsf T}L_Gx
 =\sum_{uv\in E(G)}w_{uv}(x_u-x_v)^2.
\]
For $S\subseteq V(G)$, the quadratic form counts boundary edge
weight, so
$\mathbf1_S^{\mathsf T}L_G\mathbf1_S=w(\cut_G(S))$ and
$\mathbf1_{\{v\}}^{\mathsf T}L_G\mathbf1_{\{v\}}=\deg_G(v)$.
The notation $L_G\preceq L_H$ means that
$x^{\mathsf T}L_Gx\le x^{\mathsf T}L_Hx$ for every vector $x$.
A \emph{spectral sparsifier} is a weighted subgraph with fewer
edges whose Laplacian preserves these quadratic forms within
multiplicative bounds. Finally,
$2e_G(S)=\vol_G(S)-w(\cut_G(S))$: an internal edge
contributes twice to the volume, while a crossing edge contributes
once to both the volume and the cut.

\begin{theorem}[BSS spectral sparsification
{\cite[Theorem~1.1]{BSS2012}}]
\label{thm:bss-spectral-sparsification}
For every fixed $q>1$, a loopless undirected weighted graph $G$
on $N\geq1$ vertices has a reweighted spanning subgraph $G'$ with
at most $\lceil q(N-1)\rceil$ edges such that
\[
 L_G\preceq L_{G'}\preceq
 \frac{q+1+2\sqrt q}{q+1-2\sqrt q}L_G.
\]
Such a graph can be constructed deterministically in polynomial time.
\end{theorem}

\begin{lemma}\label{lem:deterministic-density-sparsification}
Fix positive constants $C_0,\zeta,\eta$. Let $G$ be a loopless
edge-weighted graph on $N\ge1$ vertices, with parallel edge weights
already combined, of total edge weight at most one and weighted degree
at most $C_0/N$ at every vertex. Suppose that
every set of at least $\eta N$ vertices has internal edge weight
at least $\zeta$. There is a deterministic polynomial-time
construction of a simple unweighted graph $J$ with
$V(J)=V(G)$ and $E(J)\subseteq E(G)$, whose
maximum degree is bounded by a constant depending only on
$C_0,\zeta$, and whose independent sets all have size less than
$\eta N$.
\end{lemma}

\begin{proof}
Choose a fixed rational $0<\theta\le\min\{1/12,\zeta/14\}$,
and put $q=\theta^{-2}$ and $d=\lceil q\rceil$.
Apply Theorem~\ref{thm:bss-spectral-sparsification} to obtain a
reweighted subgraph $G'$ with at most
$\lceil q(N-1)\rceil\leq dN$ edges and the stated spectral bounds.
Substituting $q=\theta^{-2}$ gives
\[
 \frac{q+1+2\sqrt q}{q+1-2\sqrt q}
 =\left(\frac{1+\theta}{1-\theta}\right)^2
 =1+\frac{4\theta}{(1-\theta)^2}
 \le1+\frac{576}{121}\theta<1+5\theta,
\]
since $\theta\le1/12$.

Applying this sandwich to $\mathbf1_{\{v\}}$ gives
$\deg_G(v)\le\deg_{G'}(v)\le(1+5\theta)C_0/N$. Summing the lower
bounds over $v\in S$ gives
$\operatorname{vol}_{G'}(S)\ge\operatorname{vol}_G(S)$.
Applying the sandwich to $\mathbf1_S$ gives
$w(\cut_{G'}(S))\le(1+5\theta)w(\cut_G(S))$. Every cut has weight at most the total edge weight, so
$w(\cut_G(S))\le1$. Also $\theta\le\zeta/14$ gives
$5\theta/2\leq5\zeta/28\leq\zeta/4$. Therefore, whenever $|S|\ge\eta N$,
\begin{align*}
 2e_{G'}(S)&=\operatorname{vol}_{G'}(S)-w(\cut_{G'}(S))\\
 &\ge\operatorname{vol}_G(S)-(1+5\theta)w(\cut_G(S))\\
 &=2e_G(S)-5\theta w(\cut_G(S)),\\
 e_{G'}(S)&\ge e_G(S)-\tfrac{5\theta}{2}w(\cut_G(S))
 \ge\zeta-\tfrac{5\theta}{2}\ge3\zeta/4.
\end{align*}
Write $w'_e$ for the weight of $e\in E(G')$. Define the unweighted
graph $J$ by $V(J)=V(G)$ and
$E(J)=\{e\in E(G'):w'_e\ge\zeta/(4dN)\}$.
The edges in $E(G')\setminus E(J)$ have total weight at most
$|E(G')|\zeta/(4dN)\le\zeta/4$.
Thus every $S\subseteq V(G)$ with $|S|\ge\eta N$ satisfies
\[
 \sum_{\substack{e\in E(J)\\e\subseteq S}}w'_e
 \ge e_{G'}(S)-\sum_{e\in E(G')\setminus E(J)}w'_e
 \ge\frac{3\zeta}{4}-\frac\zeta4=\frac\zeta2>0.
\]
Consequently $J[S]$ contains an edge.
Thus every independent set of $J$ has size less than $\eta N$.

For $v\in V(J)$, every $e\in E(J)$ incident with $v$ has
$w'_e\ge\zeta/(4dN)$. Hence
\[
 \deg_J(v)\le\frac{4dN}{\zeta}\deg_{G'}(v)
 \le\frac{4(1+5\theta)C_0d}{\zeta},
\]
a constant depending only on $C_0,\zeta$. The construction is
deterministic and polynomial time by
Theorem~\ref{thm:bss-spectral-sparsification}.
\end{proof}

\begin{restatedtheorem}{Theorem~\ref{thm:bounded-degree-kr-source}}
For every fixed rational $0<\eta<1/10$, there is a constant
$d_\eta$ such that the following promise problem is Unique Games-hard.
The input is a simple unweighted graph $G$ on $n\geq1$ vertices
of maximum degree at most $d_\eta$. Distinguish the following cases:
\begin{description}
 \item[YES:] $\alphaI(G)\geq(1/2-\eta)n$.
 \item[NO:] $\alphaI(G)<\eta n$.
\end{description}
The reduction is deterministic and polynomial time.
\end{restatedtheorem}

\begin{proof}[Proof of Theorem~\ref{thm:bounded-degree-kr-source}]
Take the graph from Lemma~\ref{lem:normalized-density-source} and
apply Lemma~\ref{lem:deterministic-density-sparsification} with
$C_0=2$ and $\zeta=\zeta_\eta$. Deleting edges preserves the YES
independent set. In the NO case, every set of at least $\eta N$
vertices still contains an edge, so every independent set has size
less than $\eta N$. The output degree depends only on $\eta$.
Both constructions are deterministic and polynomial time.
\end{proof}
